% Options for packages loaded elsewhere
\PassOptionsToPackage{unicode}{hyperref}
\PassOptionsToPackage{hyphens}{url}
%
\documentclass[
]{article}
\usepackage{amsmath,amssymb}
\usepackage{iftex}
\ifPDFTeX
  \usepackage[T1]{fontenc}
  \usepackage[utf8]{inputenc}
  \usepackage{textcomp} % provide euro and other symbols
\else % if luatex or xetex
  \usepackage{unicode-math} % this also loads fontspec
  \defaultfontfeatures{Scale=MatchLowercase}
  \defaultfontfeatures[\rmfamily]{Ligatures=TeX,Scale=1}
\fi
\usepackage{lmodern}
\ifPDFTeX\else
  % xetex/luatex font selection
\fi
% Use upquote if available, for straight quotes in verbatim environments
\IfFileExists{upquote.sty}{\usepackage{upquote}}{}
\IfFileExists{microtype.sty}{% use microtype if available
  \usepackage[]{microtype}
  \UseMicrotypeSet[protrusion]{basicmath} % disable protrusion for tt fonts
}{}
\makeatletter
\@ifundefined{KOMAClassName}{% if non-KOMA class
  \IfFileExists{parskip.sty}{%
    \usepackage{parskip}
  }{% else
    \setlength{\parindent}{0pt}
    \setlength{\parskip}{6pt plus 2pt minus 1pt}}
}{% if KOMA class
  \KOMAoptions{parskip=half}}
\makeatother
\usepackage{xcolor}
\usepackage[margin=1in]{geometry}
\usepackage{color}
\usepackage{fancyvrb}
\newcommand{\VerbBar}{|}
\newcommand{\VERB}{\Verb[commandchars=\\\{\}]}
\DefineVerbatimEnvironment{Highlighting}{Verbatim}{commandchars=\\\{\}}
% Add ',fontsize=\small' for more characters per line
\usepackage{framed}
\definecolor{shadecolor}{RGB}{248,248,248}
\newenvironment{Shaded}{\begin{snugshade}}{\end{snugshade}}
\newcommand{\AlertTok}[1]{\textcolor[rgb]{0.94,0.16,0.16}{#1}}
\newcommand{\AnnotationTok}[1]{\textcolor[rgb]{0.56,0.35,0.01}{\textbf{\textit{#1}}}}
\newcommand{\AttributeTok}[1]{\textcolor[rgb]{0.13,0.29,0.53}{#1}}
\newcommand{\BaseNTok}[1]{\textcolor[rgb]{0.00,0.00,0.81}{#1}}
\newcommand{\BuiltInTok}[1]{#1}
\newcommand{\CharTok}[1]{\textcolor[rgb]{0.31,0.60,0.02}{#1}}
\newcommand{\CommentTok}[1]{\textcolor[rgb]{0.56,0.35,0.01}{\textit{#1}}}
\newcommand{\CommentVarTok}[1]{\textcolor[rgb]{0.56,0.35,0.01}{\textbf{\textit{#1}}}}
\newcommand{\ConstantTok}[1]{\textcolor[rgb]{0.56,0.35,0.01}{#1}}
\newcommand{\ControlFlowTok}[1]{\textcolor[rgb]{0.13,0.29,0.53}{\textbf{#1}}}
\newcommand{\DataTypeTok}[1]{\textcolor[rgb]{0.13,0.29,0.53}{#1}}
\newcommand{\DecValTok}[1]{\textcolor[rgb]{0.00,0.00,0.81}{#1}}
\newcommand{\DocumentationTok}[1]{\textcolor[rgb]{0.56,0.35,0.01}{\textbf{\textit{#1}}}}
\newcommand{\ErrorTok}[1]{\textcolor[rgb]{0.64,0.00,0.00}{\textbf{#1}}}
\newcommand{\ExtensionTok}[1]{#1}
\newcommand{\FloatTok}[1]{\textcolor[rgb]{0.00,0.00,0.81}{#1}}
\newcommand{\FunctionTok}[1]{\textcolor[rgb]{0.13,0.29,0.53}{\textbf{#1}}}
\newcommand{\ImportTok}[1]{#1}
\newcommand{\InformationTok}[1]{\textcolor[rgb]{0.56,0.35,0.01}{\textbf{\textit{#1}}}}
\newcommand{\KeywordTok}[1]{\textcolor[rgb]{0.13,0.29,0.53}{\textbf{#1}}}
\newcommand{\NormalTok}[1]{#1}
\newcommand{\OperatorTok}[1]{\textcolor[rgb]{0.81,0.36,0.00}{\textbf{#1}}}
\newcommand{\OtherTok}[1]{\textcolor[rgb]{0.56,0.35,0.01}{#1}}
\newcommand{\PreprocessorTok}[1]{\textcolor[rgb]{0.56,0.35,0.01}{\textit{#1}}}
\newcommand{\RegionMarkerTok}[1]{#1}
\newcommand{\SpecialCharTok}[1]{\textcolor[rgb]{0.81,0.36,0.00}{\textbf{#1}}}
\newcommand{\SpecialStringTok}[1]{\textcolor[rgb]{0.31,0.60,0.02}{#1}}
\newcommand{\StringTok}[1]{\textcolor[rgb]{0.31,0.60,0.02}{#1}}
\newcommand{\VariableTok}[1]{\textcolor[rgb]{0.00,0.00,0.00}{#1}}
\newcommand{\VerbatimStringTok}[1]{\textcolor[rgb]{0.31,0.60,0.02}{#1}}
\newcommand{\WarningTok}[1]{\textcolor[rgb]{0.56,0.35,0.01}{\textbf{\textit{#1}}}}
\usepackage{graphicx}
\makeatletter
\def\maxwidth{\ifdim\Gin@nat@width>\linewidth\linewidth\else\Gin@nat@width\fi}
\def\maxheight{\ifdim\Gin@nat@height>\textheight\textheight\else\Gin@nat@height\fi}
\makeatother
% Scale images if necessary, so that they will not overflow the page
% margins by default, and it is still possible to overwrite the defaults
% using explicit options in \includegraphics[width, height, ...]{}
\setkeys{Gin}{width=\maxwidth,height=\maxheight,keepaspectratio}
% Set default figure placement to htbp
\makeatletter
\def\fps@figure{htbp}
\makeatother
\setlength{\emergencystretch}{3em} % prevent overfull lines
\providecommand{\tightlist}{%
  \setlength{\itemsep}{0pt}\setlength{\parskip}{0pt}}
\setcounter{secnumdepth}{-\maxdimen} % remove section numbering
\ifLuaTeX
  \usepackage{selnolig}  % disable illegal ligatures
\fi
\usepackage{bookmark}
\IfFileExists{xurl.sty}{\usepackage{xurl}}{} % add URL line breaks if available
\urlstyle{same}
\hypersetup{
  pdftitle={Notes on Chapter 14: Strings},
  pdfauthor={Norman Lim},
  hidelinks,
  pdfcreator={LaTeX via pandoc}}

\title{Notes on Chapter 14: Strings}
\author{Norman Lim}
\date{2025-07-28}

\begin{document}
\maketitle

\subsubsection{Prerequisites}\label{prerequisites}

\begin{Shaded}
\begin{Highlighting}[]
\FunctionTok{library}\NormalTok{(tidyverse)}
\end{Highlighting}
\end{Shaded}

\begin{verbatim}
## -- Attaching core tidyverse packages ------------------------ tidyverse 2.0.0 --
## v dplyr     1.1.4     v readr     2.1.5
## v forcats   1.0.0     v stringr   1.5.1
## v ggplot2   3.5.2     v tibble    3.3.0
## v lubridate 1.9.4     v tidyr     1.3.1
## v purrr     1.1.0     
## -- Conflicts ------------------------------------------ tidyverse_conflicts() --
## x dplyr::filter() masks stats::filter()
## x dplyr::lag()    masks stats::lag()
## i Use the conflicted package (<http://conflicted.r-lib.org/>) to force all conflicts to become errors
\end{verbatim}

\begin{Shaded}
\begin{Highlighting}[]
\FunctionTok{library}\NormalTok{(babynames)}
\end{Highlighting}
\end{Shaded}

\subsection{Creating a string}\label{creating-a-string}

\begin{Shaded}
\begin{Highlighting}[]
\NormalTok{string1 }\OtherTok{\textless{}{-}} \StringTok{"This is a string"}
\NormalTok{string2 }\OtherTok{\textless{}{-}} \StringTok{\textquotesingle{}If I want to include a "quote" inside a string, I use single quotes\textquotesingle{}}
\end{Highlighting}
\end{Shaded}

\subsubsection{Escapes}\label{escapes}

\begin{Shaded}
\begin{Highlighting}[]
\NormalTok{double\_quote }\OtherTok{\textless{}{-}} \StringTok{"}\SpecialCharTok{\textbackslash{}"}\StringTok{"} \CommentTok{\# or \textquotesingle{}"\textquotesingle{}}
\NormalTok{single\_quote }\OtherTok{\textless{}{-}} \StringTok{\textquotesingle{}}\SpecialCharTok{\textbackslash{}\textquotesingle{}}\StringTok{\textquotesingle{}} \CommentTok{\# or "\textquotesingle{}"}
\end{Highlighting}
\end{Shaded}

\begin{itemize}
\tightlist
\item
  Using a literal backslash to escape a backlash:
\end{itemize}

\begin{Shaded}
\begin{Highlighting}[]
\NormalTok{backslash }\OtherTok{\textless{}{-}} \StringTok{"}\SpecialCharTok{\textbackslash{}\textbackslash{}}\StringTok{"}
\end{Highlighting}
\end{Shaded}

\begin{itemize}
\tightlist
\item
  To view strings properly, use \texttt{str\_view()}:
\end{itemize}

\begin{Shaded}
\begin{Highlighting}[]
\NormalTok{x }\OtherTok{\textless{}{-}} \FunctionTok{c}\NormalTok{(single\_quote, double\_quote, backslash)}

\FunctionTok{str\_view}\NormalTok{(x)}
\end{Highlighting}
\end{Shaded}

\begin{verbatim}
## [1] | '
## [2] | "
## [3] | \
\end{verbatim}

\subsubsection{Raw strings}\label{raw-strings}

\begin{itemize}
\tightlist
\item
  Using multiple quotes and/or backslashes can be hard to read and
  confusing.\\
\item
  Better method is to declare the string as ``raw'' string, just like in
  Python:
\end{itemize}

\begin{Shaded}
\begin{Highlighting}[]
\NormalTok{tricky }\OtherTok{\textless{}{-}}\NormalTok{ r}\StringTok{"(double\_quote \textless{}{-} "}\NormalTok{\textbackslash{}}\StringTok{""} \CommentTok{\# or \textquotesingle{}"\textquotesingle{}}
\NormalTok{single\_quote }\OtherTok{\textless{}{-}} \StringTok{\textquotesingle{}}\SpecialCharTok{\textbackslash{}\textquotesingle{}}\StringTok{\textquotesingle{}} \CommentTok{\# or "\textquotesingle{}")"}

\FunctionTok{str\_view}\NormalTok{(tricky)}
\end{Highlighting}
\end{Shaded}

\begin{verbatim}
## [1] | double_quote <- "\"" # or '"'
##     | single_quote <- '\'' # or "'"
\end{verbatim}

\begin{itemize}
\tightlist
\item
  Raw strings in R mus be enclosed with \texttt{r"()"} or
  \texttt{r"{[}{]}"\textasciigrave{}\textasciigrave{}\ or}r''\{\}``\,`
\end{itemize}

\subsubsection{Other special characters}\label{other-special-characters}

\begin{itemize}
\tightlist
\item
  \texttt{\textbackslash{}n} for new line.\\
\item
  \texttt{\textbackslash{}t} for tab.
\item
  \texttt{\textbackslash{}u} or \texttt{\textbackslash{}U} for unicode
  escapes
\end{itemize}

\begin{Shaded}
\begin{Highlighting}[]
\NormalTok{x }\OtherTok{\textless{}{-}} \FunctionTok{c}\NormalTok{(}\StringTok{"one}\SpecialCharTok{\textbackslash{}n}\StringTok{two"}\NormalTok{, }\StringTok{"one}\SpecialCharTok{\textbackslash{}t}\StringTok{two"}\NormalTok{, }\StringTok{"\textbackslash{}u00b5"}\NormalTok{, }\StringTok{"\textbackslash{}U0001f604"}\NormalTok{)}

\FunctionTok{str\_view}\NormalTok{(x)}
\end{Highlighting}
\end{Shaded}

\begin{verbatim}
## [1] | one
##     | two
## [2] | one{\t}two
## [3] | µ
## [4] | 😄
\end{verbatim}

\begin{itemize}
\tightlist
\item
  \texttt{str\_view()} uses `\{\t\} for tabs to make them easier to
  spot.
\end{itemize}

\subsubsection{Exercises}\label{exercises}

\begin{enumerate}
\def\labelenumi{\arabic{enumi}.}
\tightlist
\item
  Create strings that contain the following values: \textgreater{} 1. He
  said ``That's amazing!'' \textgreater{} 2. \a\b\c\d
  > 3. \textbackslash\textbackslash\textbackslash{}
\end{enumerate}

\begin{Shaded}
\begin{Highlighting}[]
\CommentTok{\# No. 1}
\NormalTok{str\_1 }\OtherTok{\textless{}{-}}\NormalTok{ r}\StringTok{"("}\NormalTok{He said That}\StringTok{\textquotesingle{}s amazing!")"}
\StringTok{str\_view(str\_1)}
\end{Highlighting}
\end{Shaded}

\begin{verbatim}
## [1] | "He said That's amazing!"
\end{verbatim}

\begin{Shaded}
\begin{Highlighting}[]
\CommentTok{\# No. 2}
\NormalTok{str\_2 }\OtherTok{\textless{}{-}}\NormalTok{ r}\StringTok{"(}\SpecialCharTok{\textbackslash{}a\textbackslash{}b}\StringTok{\textbackslash{}c\textbackslash{}d)"}
\FunctionTok{str\_view}\NormalTok{(str\_2)}
\end{Highlighting}
\end{Shaded}

\begin{verbatim}
## [1] | \a\b\c\d
\end{verbatim}

\begin{Shaded}
\begin{Highlighting}[]
\CommentTok{\# No. 3}
\NormalTok{str\_3 }\OtherTok{\textless{}{-}}\NormalTok{ r}\StringTok{"(}\SpecialCharTok{\textbackslash{}\textbackslash{}\textbackslash{}\textbackslash{}\textbackslash{}\textbackslash{}}\StringTok{)"}
\FunctionTok{str\_view}\NormalTok{(str\_3)}
\end{Highlighting}
\end{Shaded}

\begin{verbatim}
## [1] | \\\\\\
\end{verbatim}

\begin{enumerate}
\def\labelenumi{\arabic{enumi}.}
\setcounter{enumi}{1}
\tightlist
\item
  Create the string in your R session and print it. What happens to the
  special ``\u00a0''? How does \texttt{str\_view()} display it? Can you
  do a little googling to figure out what this special character is?
\end{enumerate}

Ans.

\begin{Shaded}
\begin{Highlighting}[]
\NormalTok{str\_4 }\OtherTok{\textless{}{-}} \StringTok{"This\textbackslash{}u00a0is\textbackslash{}u00a0tricky"}
\FunctionTok{str\_view}\NormalTok{(str\_4)}
\end{Highlighting}
\end{Shaded}

\begin{verbatim}
## [1] | This{\u00a0}is{\u00a0}tricky
\end{verbatim}

The unicode character u+00a0 (\texttt{\textbackslash{}u00a0}) stands for
no-break space. This is a space character that prevents an automatic
line break at its position.

\subsubsection{Creating many strings from
data}\label{creating-many-strings-from-data}

\begin{itemize}
\tightlist
\item
  \texttt{str\_c}: takes any number of vectors as agruments and returns
  a character vector:
\end{itemize}

\begin{Shaded}
\begin{Highlighting}[]
\FunctionTok{str\_c}\NormalTok{(}\StringTok{"x"}\NormalTok{, }\StringTok{"y"}\NormalTok{)}
\end{Highlighting}
\end{Shaded}

\begin{verbatim}
## [1] "xy"
\end{verbatim}

\begin{Shaded}
\begin{Highlighting}[]
\FunctionTok{str\_c}\NormalTok{(}\StringTok{"x"}\NormalTok{, }\StringTok{"y"}\NormalTok{, }\StringTok{"z"}\NormalTok{)}
\end{Highlighting}
\end{Shaded}

\begin{verbatim}
## [1] "xyz"
\end{verbatim}

\begin{Shaded}
\begin{Highlighting}[]
\FunctionTok{str\_c}\NormalTok{(}\StringTok{"Hello, "}\NormalTok{, }\FunctionTok{c}\NormalTok{(}\StringTok{"John"}\NormalTok{, }\StringTok{"Susan"}\NormalTok{))}
\end{Highlighting}
\end{Shaded}

\begin{verbatim}
## [1] "Hello, John"  "Hello, Susan"
\end{verbatim}

\begin{itemize}
\tightlist
\item
  \texttt{str\_c} is similar to \texttt{past0()}, but is designed to be
  used with \texttt{mutate()}:
\end{itemize}

\begin{Shaded}
\begin{Highlighting}[]
\NormalTok{df }\OtherTok{\textless{}{-}} \FunctionTok{tibble}\NormalTok{(}\AttributeTok{name =} \FunctionTok{c}\NormalTok{(}\StringTok{"Flora"}\NormalTok{, }\StringTok{"David"}\NormalTok{, }\StringTok{"Terra"}\NormalTok{, }\ConstantTok{NA}\NormalTok{))}
\NormalTok{df }\SpecialCharTok{|\textgreater{}} 
  \FunctionTok{mutate}\NormalTok{(}\AttributeTok{greeting =} \FunctionTok{str\_c}\NormalTok{(}\StringTok{"Hi "}\NormalTok{, name, }\StringTok{"!"}\NormalTok{))}
\end{Highlighting}
\end{Shaded}

\begin{verbatim}
## # A tibble: 4 x 2
##   name  greeting 
##   <chr> <chr>    
## 1 Flora Hi Flora!
## 2 David Hi David!
## 3 Terra Hi Terra!
## 4 <NA>  <NA>
\end{verbatim}

\begin{itemize}
\tightlist
\item
  \texttt{coalesce()} displays \texttt{NA}s in a different way:
\end{itemize}

\begin{Shaded}
\begin{Highlighting}[]
\NormalTok{df }\SpecialCharTok{|\textgreater{}} 
  \FunctionTok{mutate}\NormalTok{(}
    \AttributeTok{greeting1 =} \FunctionTok{str\_c}\NormalTok{(}\StringTok{"Hi "}\NormalTok{, }\FunctionTok{coalesce}\NormalTok{(name, }\StringTok{"you"}\NormalTok{), }\StringTok{"!"}\NormalTok{),}
    \AttributeTok{greeting2 =} \FunctionTok{coalesce}\NormalTok{(}\FunctionTok{str\_c}\NormalTok{(}\StringTok{"Hi"}\NormalTok{, name, }\StringTok{"!"}\NormalTok{))}
\NormalTok{  )}
\end{Highlighting}
\end{Shaded}

\begin{verbatim}
## # A tibble: 4 x 3
##   name  greeting1 greeting2
##   <chr> <chr>     <chr>    
## 1 Flora Hi Flora! HiFlora! 
## 2 David Hi David! HiDavid! 
## 3 Terra Hi Terra! HiTerra! 
## 4 <NA>  Hi you!   <NA>
\end{verbatim}

\subsubsection{\texorpdfstring{\texttt{str\_glue()}}{str\_glue()}}\label{str_glue}

\begin{itemize}
\tightlist
\item
  Using \texttt{str\_c()} leads to a lot of \texttt{"}, making the code
  unreadable. \texttt{str\_glue()} behaves in a similar way to f-stings
  in Python, but works in dplyr:
\end{itemize}

\begin{Shaded}
\begin{Highlighting}[]
\NormalTok{df }\SpecialCharTok{|\textgreater{}} 
  \FunctionTok{mutate}\NormalTok{(}\AttributeTok{greeting =} \FunctionTok{str\_glue}\NormalTok{(}\StringTok{"Hi \{name\}!"}\NormalTok{))}
\end{Highlighting}
\end{Shaded}

\begin{verbatim}
## # A tibble: 4 x 2
##   name  greeting 
##   <chr> <glue>   
## 1 Flora Hi Flora!
## 2 David Hi David!
## 3 Terra Hi Terra!
## 4 <NA>  Hi NA!
\end{verbatim}

\begin{itemize}
\item
  One drawback of \texttt{str\_glue()} is the way it handles
  \texttt{NA}s -- \texttt{str\_glue()} converts them to the literal
  string ``NA''.
\item
  To escape the curly brackets character ``\{'' or ``\}'' in
  \texttt{str\_glue()}, you double it:
\end{itemize}

\begin{Shaded}
\begin{Highlighting}[]
\NormalTok{df }\SpecialCharTok{|\textgreater{}} 
  \FunctionTok{mutate}\NormalTok{(}\AttributeTok{greeting =} \FunctionTok{str\_glue}\NormalTok{(}\StringTok{"\{\{Hi \{name\}!\}\}"}\NormalTok{))}
\end{Highlighting}
\end{Shaded}

\begin{verbatim}
## # A tibble: 4 x 2
##   name  greeting   
##   <chr> <glue>     
## 1 Flora {Hi Flora!}
## 2 David {Hi David!}
## 3 Terra {Hi Terra!}
## 4 <NA>  {Hi NA!}
\end{verbatim}

\subsubsection{\texorpdfstring{\texttt{str\_flatten()}}{str\_flatten()}}\label{str_flatten}

\begin{itemize}
\tightlist
\item
  \texttt{str\_flatten()} takes a character vector and combines each
  element of the vector into a single string:
\end{itemize}

\begin{Shaded}
\begin{Highlighting}[]
\FunctionTok{str\_flatten}\NormalTok{(}\FunctionTok{c}\NormalTok{(}\StringTok{"x"}\NormalTok{, }\StringTok{"y"}\NormalTok{, }\StringTok{"z"}\NormalTok{))}
\end{Highlighting}
\end{Shaded}

\begin{verbatim}
## [1] "xyz"
\end{verbatim}

\begin{Shaded}
\begin{Highlighting}[]
\FunctionTok{str\_flatten}\NormalTok{(}\FunctionTok{c}\NormalTok{(}\StringTok{"x"}\NormalTok{, }\StringTok{"y"}\NormalTok{, }\StringTok{"z"}\NormalTok{), }\StringTok{", "}\NormalTok{)}
\end{Highlighting}
\end{Shaded}

\begin{verbatim}
## [1] "x, y, z"
\end{verbatim}

\begin{Shaded}
\begin{Highlighting}[]
\FunctionTok{str\_flatten}\NormalTok{(}\FunctionTok{c}\NormalTok{(}\StringTok{"x"}\NormalTok{, }\StringTok{"y"}\NormalTok{, }\StringTok{"z"}\NormalTok{), }\StringTok{", "}\NormalTok{, }\AttributeTok{last =} \StringTok{", and "}\NormalTok{)}
\end{Highlighting}
\end{Shaded}

\begin{verbatim}
## [1] "x, y, and z"
\end{verbatim}

\begin{itemize}
\tightlist
\item
  Using \texttt{str\_flatten()} with \texttt{summarize()}:
\end{itemize}

\begin{Shaded}
\begin{Highlighting}[]
\NormalTok{df }\OtherTok{\textless{}{-}} \FunctionTok{tribble}\NormalTok{(}
  \SpecialCharTok{\textasciitilde{}}\NormalTok{ name, }\SpecialCharTok{\textasciitilde{}}\NormalTok{ fruit,}
  \StringTok{"Carmen"}\NormalTok{, }\StringTok{"banana"}\NormalTok{,}
  \StringTok{"Carmen"}\NormalTok{, }\StringTok{"apple"}\NormalTok{,}
  \StringTok{"Marvin"}\NormalTok{, }\StringTok{"nectarine"}\NormalTok{,}
  \StringTok{"Terrence"}\NormalTok{, }\StringTok{"cantaloupe"}\NormalTok{,}
  \StringTok{"Terrence"}\NormalTok{, }\StringTok{"papaya"}\NormalTok{,}
  \StringTok{"Terrence"}\NormalTok{, }\StringTok{"mandarin"}
\NormalTok{)}

\NormalTok{df }\SpecialCharTok{|\textgreater{}} 
  \FunctionTok{group\_by}\NormalTok{(name) }\SpecialCharTok{|\textgreater{}} 
  \FunctionTok{summarize}\NormalTok{(}\AttributeTok{fruits =} \FunctionTok{str\_flatten}\NormalTok{(fruit, }\StringTok{", "}\NormalTok{))}
\end{Highlighting}
\end{Shaded}

\begin{verbatim}
## # A tibble: 3 x 2
##   name     fruits                      
##   <chr>    <chr>                       
## 1 Carmen   banana, apple               
## 2 Marvin   nectarine                   
## 3 Terrence cantaloupe, papaya, mandarin
\end{verbatim}

\subsubsection{Exercises}\label{exercises-1}

\begin{enumerate}
\def\labelenumi{\arabic{enumi}.}
\tightlist
\item
  Compare and contrast the results of \texttt{paste0} with
  \texttt{str\_c()} for the following inputs:
\end{enumerate}

\begin{Shaded}
\begin{Highlighting}[]
\FunctionTok{str\_c}\NormalTok{(}\StringTok{"hi "}\NormalTok{, }\ConstantTok{NA}\NormalTok{)}
\FunctionTok{str\_c}\NormalTok{(letters[}\DecValTok{1}\SpecialCharTok{:}\DecValTok{2}\NormalTok{], letters[}\DecValTok{1}\SpecialCharTok{:}\DecValTok{3}\NormalTok{])}
\end{Highlighting}
\end{Shaded}

\begin{Shaded}
\begin{Highlighting}[]
\CommentTok{\# str\_c}
\FunctionTok{str\_c}\NormalTok{(}\StringTok{"hi "}\NormalTok{, }\ConstantTok{NA}\NormalTok{)}
\end{Highlighting}
\end{Shaded}

\begin{verbatim}
## [1] NA
\end{verbatim}

\begin{Shaded}
\begin{Highlighting}[]
\CommentTok{\# str\_c(letters[1:2], letters[1:3])}

\CommentTok{\# paste0}
\FunctionTok{paste0}\NormalTok{(}\StringTok{"hi "}\NormalTok{, }\ConstantTok{NA}\NormalTok{)}
\end{Highlighting}
\end{Shaded}

\begin{verbatim}
## [1] "hi NA"
\end{verbatim}

\begin{Shaded}
\begin{Highlighting}[]
\FunctionTok{paste0}\NormalTok{(letters[}\DecValTok{1}\SpecialCharTok{:}\DecValTok{2}\NormalTok{], letters[}\DecValTok{1}\SpecialCharTok{:}\DecValTok{3}\NormalTok{])}
\end{Highlighting}
\end{Shaded}

\begin{verbatim}
## [1] "aa" "bb" "ac"
\end{verbatim}

\texttt{str\_c()} cannot threw an error running the second line due to
``incompatible size''

\subsection{Extracting data from
strings}\label{extracting-data-from-strings}

\subsubsection{Separating into rows}\label{separating-into-rows}

\begin{itemize}
\tightlist
\item
  The most common tool is \texttt{separate\_longer\_delim()}:
\end{itemize}

\begin{Shaded}
\begin{Highlighting}[]
\NormalTok{df1 }\OtherTok{\textless{}{-}} \FunctionTok{tibble}\NormalTok{(}\AttributeTok{x =} \FunctionTok{c}\NormalTok{(}\StringTok{"a,b,c"}\NormalTok{, }\StringTok{"d,e"}\NormalTok{, }\StringTok{"f"}\NormalTok{))}
\NormalTok{df1 }\SpecialCharTok{|\textgreater{}} 
  \FunctionTok{separate\_longer\_delim}\NormalTok{(x, }\AttributeTok{delim =} \StringTok{","}\NormalTok{)}
\end{Highlighting}
\end{Shaded}

\begin{verbatim}
## # A tibble: 6 x 1
##   x    
##   <chr>
## 1 a    
## 2 b    
## 3 c    
## 4 d    
## 5 e    
## 6 f
\end{verbatim}

\begin{itemize}
\tightlist
\item
  Using \texttt{separate\_longer\_position()}:
\end{itemize}

\begin{Shaded}
\begin{Highlighting}[]
\NormalTok{df2 }\OtherTok{\textless{}{-}} \FunctionTok{tibble}\NormalTok{(}\AttributeTok{x =} \FunctionTok{c}\NormalTok{(}\StringTok{"1211"}\NormalTok{, }\StringTok{"131"}\NormalTok{, }\StringTok{"21"}\NormalTok{))}
\NormalTok{df2 }\SpecialCharTok{|\textgreater{}} 
  \FunctionTok{separate\_longer\_position}\NormalTok{(x, }\AttributeTok{width =} \DecValTok{1}\NormalTok{)}
\end{Highlighting}
\end{Shaded}

\begin{verbatim}
## # A tibble: 9 x 1
##   x    
##   <chr>
## 1 1    
## 2 2    
## 3 1    
## 4 1    
## 5 1    
## 6 3    
## 7 1    
## 8 2    
## 9 1
\end{verbatim}

\subsubsection{Separating into columns}\label{separating-into-columns}

\begin{itemize}
\tightlist
\item
  Using \texttt{separate\_wider\_delim()}:
\end{itemize}

\begin{Shaded}
\begin{Highlighting}[]
\NormalTok{df3 }\OtherTok{\textless{}{-}} \FunctionTok{tibble}\NormalTok{(}\AttributeTok{x =} \FunctionTok{c}\NormalTok{(}\StringTok{"a10.1.2022"}\NormalTok{, }\StringTok{"b10.2.2011"}\NormalTok{, }\StringTok{"e15.1.2015"}\NormalTok{))}

\NormalTok{df3 }\SpecialCharTok{|\textgreater{}} 
  \FunctionTok{separate\_wider\_delim}\NormalTok{(}
\NormalTok{    x, }\AttributeTok{delim =} \StringTok{"."}\NormalTok{,}
    \AttributeTok{names =} \FunctionTok{c}\NormalTok{(}\StringTok{"code"}\NormalTok{, }\StringTok{"edition"}\NormalTok{, }\StringTok{"year"}\NormalTok{)}
\NormalTok{  )}
\end{Highlighting}
\end{Shaded}

\begin{verbatim}
## # A tibble: 3 x 3
##   code  edition year 
##   <chr> <chr>   <chr>
## 1 a10   1       2022 
## 2 b10   2       2011 
## 3 e15   1       2015
\end{verbatim}

\begin{itemize}
\tightlist
\item
  Using the argument \texttt{NA} in \texttt{searate\_wider\_delim()} to
  omit a piece from the source vector :
\end{itemize}

\begin{Shaded}
\begin{Highlighting}[]
\NormalTok{df3 }\SpecialCharTok{|\textgreater{}} 
  \FunctionTok{separate\_wider\_delim}\NormalTok{(}
\NormalTok{    x, }
    \AttributeTok{delim =} \StringTok{"."}\NormalTok{,}
    \AttributeTok{names =} \FunctionTok{c}\NormalTok{(}\StringTok{"code"}\NormalTok{, }\ConstantTok{NA}\NormalTok{, }\StringTok{"year"}\NormalTok{)}
\NormalTok{  )}
\end{Highlighting}
\end{Shaded}

\begin{verbatim}
## # A tibble: 3 x 2
##   code  year 
##   <chr> <chr>
## 1 a10   2022 
## 2 b10   2011 
## 3 e15   2015
\end{verbatim}

\begin{itemize}
\tightlist
\item
  \texttt{separate\_wider\_position()} generally requires a known fixed
  set of columns:
\end{itemize}

\begin{Shaded}
\begin{Highlighting}[]
\NormalTok{df4 }\OtherTok{\textless{}{-}} \FunctionTok{tibble}\NormalTok{(}\AttributeTok{x =} \FunctionTok{c}\NormalTok{(}\StringTok{"202215TX"}\NormalTok{, }\StringTok{"202122LA"}\NormalTok{, }\StringTok{"202325CA"}\NormalTok{))}
\NormalTok{df4 }\SpecialCharTok{|\textgreater{}} 
  \FunctionTok{separate\_wider\_position}\NormalTok{(}
\NormalTok{    x,}
    \AttributeTok{widths =} \FunctionTok{c}\NormalTok{(}\AttributeTok{year =} \DecValTok{4}\NormalTok{, }\AttributeTok{age =} \DecValTok{2}\NormalTok{, }\AttributeTok{state =} \DecValTok{2}\NormalTok{)}
\NormalTok{  )}
\end{Highlighting}
\end{Shaded}

\begin{verbatim}
## # A tibble: 3 x 3
##   year  age   state
##   <chr> <chr> <chr>
## 1 2022  15    TX   
## 2 2021  22    LA   
## 3 2023  25    CA
\end{verbatim}

\subsubsection{Diagnosing widening
problems}\label{diagnosing-widening-problems}

-If the width of a column is not known, the tools \texttt{too\_few} and
\texttt{too\_many} can be useful:

\begin{Shaded}
\begin{Highlighting}[]
\NormalTok{df }\OtherTok{\textless{}{-}} \FunctionTok{tibble}\NormalTok{(}\AttributeTok{x =} \FunctionTok{c}\NormalTok{(}\StringTok{"1{-}1{-}1"}\NormalTok{, }\StringTok{"1{-}1{-}2"}\NormalTok{, }\StringTok{"1{-}3"}\NormalTok{, }\StringTok{"1{-}3{-}2"}\NormalTok{, }\StringTok{"1"}\NormalTok{))}
\end{Highlighting}
\end{Shaded}

Running the code below will not work:

\begin{Shaded}
\begin{Highlighting}[]
\NormalTok{df }\SpecialCharTok{|\textgreater{}} 
  \FunctionTok{separate\_wider\_delim}\NormalTok{(}
\NormalTok{    x,}
    \AttributeTok{delim =} \StringTok{"{-}"}\NormalTok{,}
    \AttributeTok{names =} \FunctionTok{c}\NormalTok{(}\StringTok{"x"}\NormalTok{, }\StringTok{"y"}\NormalTok{, }\StringTok{"z"}\NormalTok{)}
\NormalTok{  )}
\end{Highlighting}
\end{Shaded}

This code will throw the following error:

\begin{verbatim}
#> Error in `separate_wider_delim()`:
#> ! Expected 3 pieces in each element of `x`.
#> ! 2 values were too short.
#> ℹ Use `too_few = "debug"` to diagnose the problem.
#> ℹ Use `too_few = "align_start"/"align_end"` to silence this message
\end{verbatim}

\begin{itemize}
\tightlist
\item
  We can use these error messages to debug our code:
\end{itemize}

\begin{Shaded}
\begin{Highlighting}[]
\NormalTok{debug }\OtherTok{\textless{}{-}}\NormalTok{ df }\SpecialCharTok{|\textgreater{}} 
  \FunctionTok{separate\_wider\_delim}\NormalTok{(}
\NormalTok{    x,}
    \AttributeTok{delim =} \StringTok{"{-}"}\NormalTok{,}
    \AttributeTok{names =} \FunctionTok{c}\NormalTok{(}\StringTok{"x"}\NormalTok{, }\StringTok{"y"}\NormalTok{, }\StringTok{"z"}\NormalTok{),}
    \AttributeTok{too\_few =} \StringTok{"debug"}
\NormalTok{  )}
\end{Highlighting}
\end{Shaded}

\begin{verbatim}
## Warning: Debug mode activated: adding variables `x_ok`, `x_pieces`, and
## `x_remainder`.
\end{verbatim}

Looking at the contents of \texttt{debug}:

\begin{Shaded}
\begin{Highlighting}[]
\NormalTok{debug}
\end{Highlighting}
\end{Shaded}

\begin{verbatim}
## # A tibble: 5 x 6
##   x     y     z     x_ok  x_pieces x_remainder
##   <chr> <chr> <chr> <lgl>    <int> <chr>      
## 1 1-1-1 1     1     TRUE         3 ""         
## 2 1-1-2 1     2     TRUE         3 ""         
## 3 1-3   3     <NA>  FALSE        2 ""         
## 4 1-3-2 3     2     TRUE         3 ""         
## 5 1     <NA>  <NA>  FALSE        1 ""
\end{verbatim}

\begin{itemize}
\tightlist
\item
  The \texttt{x\_ok} column tells us where the problem occurred:
\end{itemize}

\begin{Shaded}
\begin{Highlighting}[]
\NormalTok{debug }\SpecialCharTok{|\textgreater{}} \FunctionTok{filter}\NormalTok{(}\SpecialCharTok{!}\NormalTok{x\_ok)}
\end{Highlighting}
\end{Shaded}

\begin{verbatim}
## # A tibble: 2 x 6
##   x     y     z     x_ok  x_pieces x_remainder
##   <chr> <chr> <chr> <lgl>    <int> <chr>      
## 1 1-3   3     <NA>  FALSE        2 ""         
## 2 1     <NA>  <NA>  FALSE        1 ""
\end{verbatim}

\begin{itemize}
\tightlist
\item
  The \texttt{x\_pieces} column how many pieces were found. In this case
  we are expecting 3, but we only found 2 and 1 pieces. Now we know that
  the problem was due to having ``too\_few'' pieces. We can fill them
  with \texttt{NA}s to avoid R from throwing an error.\\
\item
  We can then use the argument \texttt{too\_few\ =\ "align\_start"} or
  \texttt{too\_few\ =\ "align\_end"} to control where the \texttt{NA}s
  should go.
\end{itemize}

\begin{Shaded}
\begin{Highlighting}[]
\NormalTok{df }\SpecialCharTok{|\textgreater{}} 
  \FunctionTok{separate\_wider\_delim}\NormalTok{(}
\NormalTok{    x,}
    \AttributeTok{delim =} \StringTok{"{-}"}\NormalTok{,}
    \AttributeTok{names =} \FunctionTok{c}\NormalTok{(}\StringTok{"x"}\NormalTok{, }\StringTok{"y"}\NormalTok{, }\StringTok{"z"}\NormalTok{),}
    \AttributeTok{too\_few =} \StringTok{"align\_start"}
\NormalTok{  )}
\end{Highlighting}
\end{Shaded}

\begin{verbatim}
## # A tibble: 5 x 3
##   x     y     z    
##   <chr> <chr> <chr>
## 1 1     1     1    
## 2 1     1     2    
## 3 1     3     <NA> 
## 4 1     3     2    
## 5 1     <NA>  <NA>
\end{verbatim}

What if we have too many pieces?

\begin{Shaded}
\begin{Highlighting}[]
\NormalTok{df }\OtherTok{\textless{}{-}} \FunctionTok{tibble}\NormalTok{(}\AttributeTok{x =} \FunctionTok{c}\NormalTok{(}\StringTok{"1{-}1{-}1"}\NormalTok{, }\StringTok{"1{-}1{-}2"}\NormalTok{, }\StringTok{"1{-}3{-}5{-}6"}\NormalTok{, }\StringTok{"1{-}3{-}2"}\NormalTok{, }\StringTok{"1{-}3{-}5{-}7{-}9"}\NormalTok{))}
\end{Highlighting}
\end{Shaded}

\begin{itemize}
\tightlist
\item
  Using the argument \texttt{too\_many}:
\end{itemize}

\begin{Shaded}
\begin{Highlighting}[]
\NormalTok{debug }\OtherTok{\textless{}{-}}\NormalTok{ df }\SpecialCharTok{|\textgreater{}} 
  \FunctionTok{separate\_wider\_delim}\NormalTok{(}
\NormalTok{    x,}
    \AttributeTok{delim =} \StringTok{"{-}"}\NormalTok{,}
    \AttributeTok{names =} \FunctionTok{c}\NormalTok{(}\StringTok{"x"}\NormalTok{, }\StringTok{"y"}\NormalTok{, }\StringTok{"z"}\NormalTok{),}
    \AttributeTok{too\_many =} \StringTok{"debug"}
\NormalTok{  )}
\end{Highlighting}
\end{Shaded}

\begin{verbatim}
## Warning: Debug mode activated: adding variables `x_ok`, `x_pieces`, and
## `x_remainder`.
\end{verbatim}

\begin{Shaded}
\begin{Highlighting}[]
\NormalTok{debug }\SpecialCharTok{|\textgreater{}} \FunctionTok{filter}\NormalTok{(}\SpecialCharTok{!}\NormalTok{x\_ok)}
\end{Highlighting}
\end{Shaded}

\begin{verbatim}
## # A tibble: 2 x 6
##   x         y     z     x_ok  x_pieces x_remainder
##   <chr>     <chr> <chr> <lgl>    <int> <chr>      
## 1 1-3-5-6   3     5     FALSE        4 -6         
## 2 1-3-5-7-9 3     5     FALSE        5 -7-9
\end{verbatim}

\begin{itemize}
\tightlist
\item
  The \texttt{x\_remainder} column also shows the excess pieces.\\
\item
  We can opt to drop the excess pieces or merge them into the last
  column:
\end{itemize}

\begin{Shaded}
\begin{Highlighting}[]
\CommentTok{\# dropping the excess pieces}
\NormalTok{df }\SpecialCharTok{|\textgreater{}} 
  \FunctionTok{separate\_wider\_delim}\NormalTok{(}
\NormalTok{    x, }
    \AttributeTok{delim =} \StringTok{"{-}"}\NormalTok{,}
    \AttributeTok{names =} \FunctionTok{c}\NormalTok{(}\StringTok{"x"}\NormalTok{, }\StringTok{"y"}\NormalTok{, }\StringTok{"z"}\NormalTok{),}
    \AttributeTok{too\_many =} \StringTok{"drop"}
\NormalTok{  )}
\end{Highlighting}
\end{Shaded}

\begin{verbatim}
## # A tibble: 5 x 3
##   x     y     z    
##   <chr> <chr> <chr>
## 1 1     1     1    
## 2 1     1     2    
## 3 1     3     5    
## 4 1     3     2    
## 5 1     3     5
\end{verbatim}

\begin{Shaded}
\begin{Highlighting}[]
\CommentTok{\# merging the excess into the last column}
\NormalTok{df }\SpecialCharTok{|\textgreater{}} 
  \FunctionTok{separate\_wider\_delim}\NormalTok{(}
\NormalTok{    x,}
    \AttributeTok{delim =} \StringTok{"{-}"}\NormalTok{,}
    \AttributeTok{names =} \FunctionTok{c}\NormalTok{(}\StringTok{"x"}\NormalTok{, }\StringTok{"y"}\NormalTok{, }\StringTok{"z"}\NormalTok{),}
    \AttributeTok{too\_many =} \StringTok{"merge"}
\NormalTok{  )}
\end{Highlighting}
\end{Shaded}

\begin{verbatim}
## # A tibble: 5 x 3
##   x     y     z    
##   <chr> <chr> <chr>
## 1 1     1     1    
## 2 1     1     2    
## 3 1     3     5-6  
## 4 1     3     2    
## 5 1     3     5-7-9
\end{verbatim}

\subsection{Working with individual letters within a
string}\label{working-with-individual-letters-within-a-string}

\subsubsection{Length}\label{length}

\begin{itemize}
\tightlist
\item
  Getting the string length using \texttt{str\_length}:
\end{itemize}

\begin{Shaded}
\begin{Highlighting}[]
\FunctionTok{str\_length}\NormalTok{(}\FunctionTok{c}\NormalTok{(}\StringTok{"a"}\NormalTok{, }\StringTok{"R for data science"}\NormalTok{, }\ConstantTok{NA}\NormalTok{))}
\end{Highlighting}
\end{Shaded}

\begin{verbatim}
## [1]  1 18 NA
\end{verbatim}

\begin{itemize}
\tightlist
\item
  Using \texttt{str\_length()} in dplyr:
\end{itemize}

\begin{Shaded}
\begin{Highlighting}[]
\NormalTok{babynames }\SpecialCharTok{|\textgreater{}} 
  \FunctionTok{count}\NormalTok{(}\AttributeTok{length =} \FunctionTok{str\_length}\NormalTok{(name), }\AttributeTok{wt =}\NormalTok{ n)}
\end{Highlighting}
\end{Shaded}

\begin{verbatim}
## # A tibble: 14 x 2
##    length        n
##     <int>    <int>
##  1      2   338150
##  2      3  8589596
##  3      4 48506739
##  4      5 87011607
##  5      6 90749404
##  6      7 72120767
##  7      8 25404066
##  8      9 11926551
##  9     10  1306159
## 10     11  2135827
## 11     12    16295
## 12     13    10845
## 13     14     3681
## 14     15      830
\end{verbatim}

\begin{itemize}
\tightlist
\item
  Another example:
\end{itemize}

\begin{Shaded}
\begin{Highlighting}[]
\NormalTok{babynames }\SpecialCharTok{|\textgreater{}} 
  \FunctionTok{filter}\NormalTok{(}\FunctionTok{str\_length}\NormalTok{(name) }\SpecialCharTok{==} \DecValTok{15}\NormalTok{) }\SpecialCharTok{|\textgreater{}} 
  \FunctionTok{count}\NormalTok{(name, }\AttributeTok{wt =}\NormalTok{ n, }\AttributeTok{sort =} \ConstantTok{TRUE}\NormalTok{)}
\end{Highlighting}
\end{Shaded}

\begin{verbatim}
## # A tibble: 34 x 2
##    name                n
##    <chr>           <int>
##  1 Franciscojavier   123
##  2 Christopherjohn   118
##  3 Johnchristopher   118
##  4 Christopherjame   108
##  5 Christophermich    52
##  6 Ryanchristopher    45
##  7 Mariadelosangel    28
##  8 Jonathanmichael    25
##  9 Christianjoseph    22
## 10 Christopherjose    22
## # i 24 more rows
\end{verbatim}

\subsubsection{Subsetting}\label{subsetting}

\begin{itemize}
\tightlist
\item
  Using \texttt{str\_sub()}:
\end{itemize}

\begin{Shaded}
\begin{Highlighting}[]
\NormalTok{x }\OtherTok{\textless{}{-}} \FunctionTok{c}\NormalTok{(}\StringTok{"Apple"}\NormalTok{, }\StringTok{"Banana"}\NormalTok{, }\StringTok{"Pear"}\NormalTok{)}

\FunctionTok{str\_sub}\NormalTok{(x, }\DecValTok{1}\NormalTok{, }\DecValTok{3}\NormalTok{)}
\end{Highlighting}
\end{Shaded}

\begin{verbatim}
## [1] "App" "Ban" "Pea"
\end{verbatim}

\begin{itemize}
\item
  The length of the subset is equal to \texttt{end\ -\ start\ +\ 1}.
\item
  Similar to Python, we can use \texttt{-1} to refer to the last
  character, -2 to the second to last character, etc. (but this is not
  0-indexed!)
\end{itemize}

\begin{Shaded}
\begin{Highlighting}[]
\FunctionTok{str\_sub}\NormalTok{(x, }\SpecialCharTok{{-}}\DecValTok{3}\NormalTok{, }\SpecialCharTok{{-}}\DecValTok{1}\NormalTok{)}
\end{Highlighting}
\end{Shaded}

\begin{verbatim}
## [1] "ple" "ana" "ear"
\end{verbatim}

\begin{itemize}
\tightlist
\item
  But unlike Python, subsetting won't fail here if the string is too
  short -- it will just return as much as possible.
\end{itemize}

\begin{Shaded}
\begin{Highlighting}[]
\FunctionTok{str\_sub}\NormalTok{(}\StringTok{"a"}\NormalTok{, }\DecValTok{1}\NormalTok{, }\DecValTok{5}\NormalTok{)}
\end{Highlighting}
\end{Shaded}

\begin{verbatim}
## [1] "a"
\end{verbatim}

\begin{itemize}
\tightlist
\item
  Using \texttt{str\_sub()} in dplyr:
\end{itemize}

\begin{Shaded}
\begin{Highlighting}[]
\NormalTok{babynames }\SpecialCharTok{|\textgreater{}} 
  \FunctionTok{mutate}\NormalTok{(}
    \AttributeTok{first =} \FunctionTok{str\_sub}\NormalTok{(name, }\DecValTok{1}\NormalTok{, }\DecValTok{1}\NormalTok{),}
    \AttributeTok{last =} \FunctionTok{str\_sub}\NormalTok{(name, }\SpecialCharTok{{-}}\DecValTok{1}\NormalTok{, }\SpecialCharTok{{-}}\DecValTok{1}\NormalTok{)}
\NormalTok{  )}
\end{Highlighting}
\end{Shaded}

\begin{verbatim}
## # A tibble: 1,924,665 x 7
##     year sex   name          n   prop first last 
##    <dbl> <chr> <chr>     <int>  <dbl> <chr> <chr>
##  1  1880 F     Mary       7065 0.0724 M     y    
##  2  1880 F     Anna       2604 0.0267 A     a    
##  3  1880 F     Emma       2003 0.0205 E     a    
##  4  1880 F     Elizabeth  1939 0.0199 E     h    
##  5  1880 F     Minnie     1746 0.0179 M     e    
##  6  1880 F     Margaret   1578 0.0162 M     t    
##  7  1880 F     Ida        1472 0.0151 I     a    
##  8  1880 F     Alice      1414 0.0145 A     e    
##  9  1880 F     Bertha     1320 0.0135 B     a    
## 10  1880 F     Sarah      1288 0.0132 S     h    
## # i 1,924,655 more rows
\end{verbatim}

\subsection{Non-English text}\label{non-english-text}

\begin{itemize}
\tightlist
\item
  Using \texttt{charToRaw()} to convert a character to its hexadecimal
  form (aka the ASCII encoding):
\end{itemize}

\begin{Shaded}
\begin{Highlighting}[]
\FunctionTok{charToRaw}\NormalTok{(}\StringTok{"Hadley"}\NormalTok{)}
\end{Highlighting}
\end{Shaded}

\begin{verbatim}
## [1] 48 61 64 6c 65 79
\end{verbatim}

\begin{itemize}
\item
  While ASCII coverage for symbols is limited, UTF-8 can encode just
  about every character, as well as other symbols, like smileys and
  other icons and shapes.
\item
  UTF-8 can also ``fail'' when the input uses a different character
  encoding. Here is an example using \texttt{read\_csv} from the
  \textbf{readr} package:
\end{itemize}

\begin{Shaded}
\begin{Highlighting}[]
\NormalTok{x1 }\OtherTok{\textless{}{-}} \StringTok{"text}\SpecialCharTok{\textbackslash{}n}\StringTok{El Ni}\SpecialCharTok{\textbackslash{}xf1}\StringTok{o was particularly bad this year"}
\FunctionTok{read\_csv}\NormalTok{(x1)}\SpecialCharTok{$}\NormalTok{text}

\NormalTok{x2 }\OtherTok{\textless{}{-}} \StringTok{"text}\SpecialCharTok{\textbackslash{}n\textbackslash{}x82\textbackslash{}xb1\textbackslash{}x82\textbackslash{}xf1\textbackslash{}x82\textbackslash{}xc9\textbackslash{}x82\textbackslash{}xbf\textbackslash{}x82\textbackslash{}xcd}\StringTok{"}
\FunctionTok{read\_csv}\NormalTok{(x2)}\SpecialCharTok{$}\NormalTok{text}
\end{Highlighting}
\end{Shaded}

\begin{itemize}
\tightlist
\item
  We can specify the correct encoding via the \texttt{locale} argument.
\end{itemize}

\begin{Shaded}
\begin{Highlighting}[]
\FunctionTok{read\_csv}\NormalTok{(x1, }\AttributeTok{locale =} \FunctionTok{locale}\NormalTok{(}\AttributeTok{encoding =} \StringTok{"Latin1"}\NormalTok{))}\SpecialCharTok{$}\NormalTok{text}
\FunctionTok{read\_csv}\NormalTok{(x2, }\AttributeTok{locale =} \FunctionTok{locale}\NormalTok{(}\AttributeTok{encoding =} \StringTok{"Shift{-}JIS"}\NormalTok{))}\SpecialCharTok{$}\NormalTok{text}
\end{Highlighting}
\end{Shaded}

\subsubsection{Letter variations}\label{letter-variations}

\begin{itemize}
\tightlist
\item
  These two letters might look the same:
\end{itemize}

\begin{Shaded}
\begin{Highlighting}[]
\NormalTok{u }\OtherTok{\textless{}{-}} \FunctionTok{c}\NormalTok{(}\StringTok{"\textbackslash{}u00fc"}\NormalTok{, }\StringTok{"u\textbackslash{}u0308"}\NormalTok{)}
\FunctionTok{str\_view}\NormalTok{(u)}
\end{Highlighting}
\end{Shaded}

\begin{verbatim}
## [1] | ü
## [2] | ü
\end{verbatim}

\begin{itemize}
\tightlist
\item
  But they differ in length:
\end{itemize}

\begin{Shaded}
\begin{Highlighting}[]
\FunctionTok{str\_length}\NormalTok{(u)}
\end{Highlighting}
\end{Shaded}

\begin{verbatim}
## [1] 1 2
\end{verbatim}

\begin{itemize}
\tightlist
\item
  However, the string comparator function \texttt{str\_equal()} from
  \textbf{stringr} still recognizes that both have the same appearance:
\end{itemize}

\begin{Shaded}
\begin{Highlighting}[]
\NormalTok{u[[}\DecValTok{1}\NormalTok{]] }\SpecialCharTok{==}\NormalTok{ u[[}\DecValTok{2}\NormalTok{]]}
\end{Highlighting}
\end{Shaded}

\begin{verbatim}
## [1] FALSE
\end{verbatim}

\begin{Shaded}
\begin{Highlighting}[]
\FunctionTok{str\_equal}\NormalTok{(u[[}\DecValTok{1}\NormalTok{]], u[[}\DecValTok{2}\NormalTok{]])}
\end{Highlighting}
\end{Shaded}

\begin{verbatim}
## [1] TRUE
\end{verbatim}

\subsubsection{Locale-dependent
functions}\label{locale-dependent-functions}

\begin{itemize}
\item
  To check which locale is supported in \textbf{stringr}, use this
  command: \texttt{stringi::stri\_locale\_list()}.
\item
  Example of difference in language rules depending on locale:
\end{itemize}

\begin{Shaded}
\begin{Highlighting}[]
\FunctionTok{str\_to\_upper}\NormalTok{(}\FunctionTok{c}\NormalTok{(}\StringTok{"i"}\NormalTok{, }\StringTok{"ı"}\NormalTok{))}
\end{Highlighting}
\end{Shaded}

\begin{verbatim}
## [1] "I" "I"
\end{verbatim}

\begin{Shaded}
\begin{Highlighting}[]
\FunctionTok{str\_to\_upper}\NormalTok{(}\FunctionTok{c}\NormalTok{(}\StringTok{"i"}\NormalTok{, }\StringTok{"ı"}\NormalTok{), }\AttributeTok{locale =} \StringTok{"tr"}\NormalTok{)}
\end{Highlighting}
\end{Shaded}

\begin{verbatim}
## [1] "İ" "I"
\end{verbatim}

\begin{itemize}
\tightlist
\item
  Example of difference in sorting rules depending on locale:
\end{itemize}

\begin{Shaded}
\begin{Highlighting}[]
\FunctionTok{str\_sort}\NormalTok{(}\FunctionTok{c}\NormalTok{(}\StringTok{"a"}\NormalTok{, }\StringTok{"c"}\NormalTok{, }\StringTok{"ch"}\NormalTok{, }\StringTok{"h"}\NormalTok{, }\StringTok{"z"}\NormalTok{))}
\end{Highlighting}
\end{Shaded}

\begin{verbatim}
## [1] "a"  "c"  "ch" "h"  "z"
\end{verbatim}

\begin{Shaded}
\begin{Highlighting}[]
\FunctionTok{str\_sort}\NormalTok{(}\FunctionTok{c}\NormalTok{(}\StringTok{"a"}\NormalTok{, }\StringTok{"c"}\NormalTok{, }\StringTok{"ch"}\NormalTok{, }\StringTok{"h"}\NormalTok{, }\StringTok{"z"}\NormalTok{), }\AttributeTok{locale =} \StringTok{"cs"}\NormalTok{)}
\end{Highlighting}
\end{Shaded}

\begin{verbatim}
## [1] "a"  "c"  "h"  "ch" "z"
\end{verbatim}

\begin{itemize}
\tightlist
\item
  The \texttt{arrange()} function in \textbf{dplyr} is also
  locale-dependent.
\end{itemize}

\end{document}
